% !TeX root = ../sections/02_exercises.tex
\subsection{2_2}
\begin{exercise}[情報数学 A 中間試験 2]
	\(-\pi \leq x \leq \pi\) において，関数
	\[
	f(x)=|x|
	\]
	を
	\[
	T_n(x)=\frac{\alpha_0}{2}
	+\sum_{k=1}^{n}
	\left(
	\alpha_k\cos kx+\beta_k\sin kx
	\right)
	\]
	により近似する．
	
	\(n=30\) としたときに，\(f(x)\) と \(T_n(x)\) の平均 \(2\) 乗誤差を
	最小にする \(\alpha_1\) と \(\beta_1\) の値をそれぞれ答えよ．
	
	\medskip
	
	※解答の際に，計算の過程は不要（求める方法は自由）で，
	答の数値のみ記せばよい．
\end{exercise}
\begin{background}
	この問題では，関数
	\[
	f(x)=|x|
	\]
	を三角多項式
	\[
	T_n(x)
	=
	\frac{\alpha_0}{2}
	+
	\sum_{k=1}^{n}
	\left(
	\alpha_k\cos kx+\beta_k\sin kx
	\right)
	\]
	によって近似する。
	
	ここで問われているのは，平均 \(2\) 乗誤差を最小にする係数である。
	
	平均 \(2\) 乗誤差とは，近似の誤差
	\[
	f(x)-T_n(x)
	\]
	の二乗を積分して平均したものである。
	典型的には，
	\[
	\frac{1}{2\pi}
	\int_{-\pi}^{\pi}
	\left|f(x)-T_n(x)\right|^2\,dx
	\]
	を考える。
	
	この量を最小にする三角多項式は，フーリエ部分和である。
	つまり，最良近似を与える係数は，\(f(x)\) のフーリエ係数である。
	
	したがって，
	\[
	\alpha_k=a_k,\qquad \beta_k=b_k
	\]
	となる。
	
	この背景には，三角関数系
	\[
	1,\quad \cos x,\quad \sin x,\quad \cos 2x,\quad \sin 2x,\quad \ldots
	\]
	が内積
	\[
	(f,g)=\int_{-\pi}^{\pi}f(x)g(x)\,dx
	\]
	に関して互いに直交する，という事実がある。
	
	平均 \(2\) 乗誤差を最小にする近似とは，
	関数 \(f(x)\) を三角関数で張られる有限次元部分空間へ直交射影することに対応する。
	その直交射影の係数がフーリエ係数である。
	
	この問題では，\(n=30\) と書かれているが，
	求めるのは \(\alpha_1\) と \(\beta_1\) だけである。
	したがって，第1フーリエ係数だけを見ればよい。
\end{background}
\begin{strategy}
	この問題を解くときは，まず
	\[
	\text{平均 \(2\) 乗誤差を最小にする係数}
	=
	\text{フーリエ係数}
	\]
	であることを思い出す。
	
	したがって，
	\[
	\alpha_1=a_1,\qquad \beta_1=b_1
	\]
	を求めればよい。
	
	次に，
	\[
	f(x)=|x|
	\]
	が偶関数であることに注目する。
	偶関数に対しては，正弦係数はすべて \(0\) になる。
	したがって，
	\[
	\beta_1=b_1=0
	\]
	である。
	
	残るのは
	\[
	\alpha_1=a_1
	\]
	である。
	
	フーリエ係数の公式より，
	\[
	a_1=\frac{1}{\pi}\int_{-\pi}^{\pi}|x|\cos x\,dx
	\]
	である。
	
	\(|x|\cos x\) は偶関数なので，
	\[
	a_1=\frac{2}{\pi}\int_0^\pi x\cos x\,dx
	\]
	とできる。
	
	この積分は部分積分で計算する。
	\[
	u=x,\qquad dv=\cos x\,dx
	\]
	とおくと，
	\[
	du=dx,\qquad v=\sin x
	\]
	である。
	
	したがって，
	\[
	\int_0^\pi x\cos x\,dx
	=
	[x\sin x]_0^\pi-\int_0^\pi \sin x\,dx
	\]
	である。
	
	ここで
	\[
	[x\sin x]_0^\pi=0
	\]
	かつ
	\[
	\int_0^\pi \sin x\,dx=2
	\]
	なので，
	\[
	\int_0^\pi x\cos x\,dx=-2
	\]
	である。
	
	よって，
	\[
	a_1=\frac{2}{\pi}(-2)=-\frac{4}{\pi}
	\]
	となる。
	
	したがって，
	\[
	\alpha_1=-\frac{4}{\pi},
	\qquad
	\beta_1=0
	\]
	である。
	
	この問題の発想は，
	\[
	\text{平均 \(2\) 乗誤差最小}
	\Rightarrow
	\text{フーリエ係数}
	\Rightarrow
	|x|\text{ は偶関数}
	\Rightarrow
	\beta_1=0,\ \alpha_1=a_1
	\]
	という流れである。
	
	なお，\(n=30\) という条件は，
	近似に \(\cos 30x,\sin 30x\) まで使えることを意味するが，
	\(\alpha_1,\beta_1\) の値そのものには影響しない。
\end{strategy}